% =====================================================================
%  chunks/ch03g.tex —— 原书 p.91–102（PDF 98–109）
%  第 3 章 §6（续）：Steenrod 代数、Milnor 运算、Massey 括号、Adams 谱序列
% =====================================================================
% [原文 p.91]
\begin{enumerate}[label=(\roman*),leftmargin=2.2em,start=3]
\item $Sq^j(z_1\cdot z_2)=\sum\limits_{k+l=j}Sq^k(z_1)\cdot Sq^l(z_2)$。
\item 当 $j<0$ 时 $Sq^j=0$；若 $z$ 的维数小于 $n$，则 $Sq^n(z)=0$。
\item Steenrod 平方是可加的（与 Pontryagin 幂不同）：
\[
Sq^j(z_1+z_2)=Sq^j(z_1)+Sq^j(z_2).
\]
\end{enumerate}

Steenrod 幂具有类似的性质：
\begin{enumerate}[label=(\roman*),leftmargin=2.2em]
\item $St^0_p=1$，$St^1_p=\beta$；当 $j<0$ 或 $j\not\equiv 0,1\pmod{2p-2}$ 时 $St^j_p=0$。
\item Steenrod 幂是\emph{稳定的}，即它们与悬挂同构交换。
\item 记 $\mathcal{P}^j=St^{2j(p-1)}_p$，则
\begin{equation}
St^{2j(p-1)+1}_p=\beta\mathcal{P}^j, \tag{6.9}
\end{equation}
并且还有下述公式：
\begin{equation}
\begin{aligned}
\mathcal{P}^j(z_1\cdot z_2)&=\sum_{i+k=j}\mathcal{P}^i(z_1)\cdot\mathcal{P}^k(z_2),\\
\beta(z_1\cdot z_2)&=\beta(z_1)\cdot z_2+(-1)^m z_1\cdot\beta(z_2),
\end{aligned}
\tag{6.10}
\end{equation}
其中 $m=\dim z_1$。
\item 对 $z\in H^{2n}(X;\Z/p)$ 有 $\mathcal{P}^n(z)=z^p$；而当 $j>n$ 时 $\mathcal{P}^j(z)=0$。
\item Steenrod 幂是可加的：
\[
St^j_p(z_1+z_2)=St^j_p(z_1)+St^j_p(z_2).
\]
\end{enumerate}

由 Serre（$p=2$ 的情形）与 Cartan（$p>2$ 的情形）20 世纪 50 年代初的结果，全部稳定运算构成的 $\Z/p$-代数（即 \emph{Steenrod 代数} $\mathcal{A}_p$）由 Steenrod 运算（$p=2$ 时是诸 $Sq^i$，$p$ 为奇素数时是诸 $St^i_p$）乘性地生成；由此可知，所有这些上同调运算都是可加的。

记 $\mathcal{A}^n_p$ 为 Steenrod 代数 $\mathcal{A}_p$ 中由次数为 $n$ 的运算组成的 $\Z/p$-子空间，即 $a\in\mathcal{A}^n_p$ 当且仅当 $a:H^q\to H^{q+n}$。当 $p=2$ 时，上同调运算 $a\in\mathcal{A}^n_2$ 是零运算，当且仅当元素 $a(u)$ 对适当的基本类 $u$ 为零：
\[
u=u_1\cdots u_N\in H^N(\RP^\infty_1\times\cdots\times\RP^\infty_N;\Z/2),\qquad u_j\in H^1(\RP^\infty_j;\Z/2),
\]

% [原文 p.92]
其中 $N>n$。（回忆 $\RP^\infty=K(\Z/2;1)$，且环 $H^*(\RP^\infty_1\times\cdots\times\RP^\infty_N;\Z/2)$ 是多项式环。）由上面列出的 Steenrod 平方的性质可以推知，像 $\mathcal{A}_2(u)\subset H^*(\RP^\infty_1\times\cdots\times\RP^\infty_N;\Z/2)$ 是由变量 $u_1,\dots,u_N$ 的、形状如下的那些对称多项式 $u_\omega$ 张成的 $\Z/2$-向量空间：
\[
u_\omega=\Bigl(\sum_{i_1\ne i_2\ne\cdots\ne i_s}u^{\omega_1}_{i_1}\cdots u^{\omega_s}_{i_s}\Bigr)u,\qquad \sum_j\omega_j<n,
\]
\[
\omega_j=2^{h_j}-1,\qquad \omega=(\omega_1,\dots,\omega_k),\qquad \omega_j\ge 0.
\]
把由
\begin{equation}
Sq^\omega(u_1\cdots u_s)=u_\omega,\qquad \deg Sq^\omega=\sum_j\omega_j \tag{6.11}
\end{equation}
定义的运算记作 $Sq^\omega$。注意
\[
Sq^i=Sq^\omega,\qquad\text{其中 } \omega=(\underbrace{1,\dots,1}_{i\text{ 个}}).
\]
成立下述乘积公式：
\[
Sq^\omega(z_1z_2)=\sum_{(\omega_1,\omega_2)=\omega}Sq^{\omega_1}(z_1)Sq^{\omega_2}(z_2).
\]

% [Tutor Note: 这里把 $\mathcal{A}_2$ 在 $H^*(\RP^\infty\times\cdots)$ 上的像写成「对称多项式张成」，实质是借用对称群作用下的不变量理论来编码 Steenrod 平方的乘性结构；读不懂形式定义时，先算 $N=1$ 的情形即可看出 $\mathcal{A}_2$ 由 $Sq^{2^i}$ 生成。]

对 $p>2$，考虑无穷维透镜空间 $K(\Z/p;1)=S^\infty/\Z/p$ 的若干个拷贝 $X_j$。此时上同调环 $H^*(X_j;\Z/p)$ 有生成元
\[
v_j\in H^1(X_j;\Z/p),\qquad u_j\in H^2(X_j;\Z/p),\qquad u_j=\beta v_j.
\]
在这种情形下，运算 $a\in\mathcal{A}^n_p$（$a:H^q\to H^{q+n}$）为零运算，当且仅当 $a(v)=0$，其中
\[
v=v_1\cdots v_{N_1}u_{N_1+1}\cdots u_{N_1+N_2}\in H^{N_1+N_2}(X_1\times\cdots\times X_{N_1+N_2};\Z/p),
\]
而 $N_1,N_2>n$。像 $\mathcal{A}_p(v)$ 由形状如下的多项式张成：
\[
u_{\omega',\omega''}=u\cdot\!\!\sum_{\substack{i_1<\cdots<i_s\le N_1\\ N_2\ge j_1>\cdots>j_q>N_1}}\!\!\bigl(u^{\omega'_1+1}_{i_1}\cdots u^{\omega'_s+1}_{i_s}u^{\omega''_1+1}_{j_1}\cdots u^{\omega''_q+1}_{j_q}\bigr)/v_{i_1}\cdots v_{i_s},
\]
\[
\omega'_1<\cdots<\omega'_s,\quad 1+\omega'_i=p^{h_i},\qquad \omega''_1<\cdots<\omega''_q,\quad 1+\omega''_j=p^{l_j},\qquad h_i\ge 0,\ l_j\ge 0,
\]
\[
\deg a=\sum_i(2\omega'_i+1)+\sum_j(2\omega''_j+1)=2\Bigl(\sum_i\omega'_i+\sum_j\omega''_j\Bigr)+s,
\]
\[
a:H^q(X;\Z/p)\longrightarrow H^{q+\deg a}(X;\Z/p),\qquad a=S_{(\omega',\omega'')}.
\]

% [原文 p.93]
此公式中出现的数 $s\ge 0$ 称为相应运算 $a$ 的 \emph{$\tau$-Cartan 型}，通常记作 $s(a)$ 或 $\tau(a)$。型 $s$ 与次数 $\deg a$（$s\le\deg a$）给出代数 $\mathcal{A}_p$ 的一个双分次（$p\ge 2$）：
\[
a\longmapsto[s(a),\deg a-s(a)].
\]
特别令人感兴趣的是运算 $a=S_{k,\omega}$，此时自然有
\[
k+2\sum\omega_j<N_1+2N_2.
\]
运算 $S_{k,0}=\theta_k$ 称为 \emph{Milnor 运算}。它们定义一个外代数：
\[
\theta_k\theta_j=-\theta_j\theta_k,\qquad s(\theta_k)=1,
\]
并具有进一步的性质
\[
\theta_k(z_1\cdot z_2)=\theta_k(z_1)z_2\pm z_1\theta_k(z_2),
\]
这是一般公式
\[
S_{(\omega',\omega'')}(z_1\cdot z_2)=\sum_{\substack{(\omega'_1,\omega'_2)=\omega'\\ (\omega''_1,\omega''_2)=\omega''}}(\pm 1)S_{(\omega'_1,\omega''_1)}(z_1)S_{(\omega'_2,\omega''_2)}(z_2)
\]
的特例。

% [Tutor Note: $\tau$-Cartan 型 $s(a)$ 是 Cartan 用来把 Steenrod 代数做成双分次代数的第二个分次；Milnor 后来把这个双分次代数显式写成了多项式代数的对偶，这就是下文 Adams 谱序列 $E_2$ 项的来源，读时应把它与「次数」$\deg a$ 区分开。]

% [原文 p.94]
结果表明，全部（不必稳定的）$\Z/p$-上同调运算构成的代数，由 Steenrod 运算与 Bockstein 算子 $\beta$ 通过复合（连同上同调环中的乘积运算）生成，而这些运算之间的全部关系都可由上面列出的性质推出。对其它系数群，则例如有 Pontryagin 幂；全部上同调运算的分类是 Cartan 在 20 世纪 50 年代中叶给出的。在此前的几年中，Serre 已证明：特征为 0 的域 $k$（例如 $\Q$、$\R$、$\C$）上的上同调运算全是平凡的，即本质上化归为上同调乘积运算。在这类情形下，除乘以标量 $z\mapsto\lambda z$（$\lambda\in k$）之外，没有稳定运算。

作为单值、处处有定义的上同调运算的推广，人们也考虑多值的、部分有定义的运算，如下面两个例子所示。

\begin{example}{1}
考虑 Bockstein 同态
\[
\beta:H^k(X;\Z/p)\longrightarrow H^{k+1}(X;\Z/p).
\]
如上面 §4 所述，$\beta$ 最初定义在整上链 $\tilde z\in C^k(X;\Z)$ 上，这些上链模 $p$ 约化为 $\Z/p$ 上的上闭链 $z$；然后 $\beta(z)$ 定义为 $\frac{1}{p}\partial^*\tilde z$ 模 $p$ 约化。若 $z\in\Ker\beta$，即上闭链 $\frac{1}{p}\partial^*\tilde z$（模 $p$）上同调于 0，则 $\frac{1}{p}\partial^*\tilde z=pu+\partial^*y$，于是可以定义
\[
\beta_2(z)=\frac{1}{p^2}\partial^*\tilde z\ (\mathrm{mod}\ p),
\]
只要我们在模 $\Image\beta$ 的意义下工作。归纳地，在 $\Ker\beta_{k-1}$（$k\ge 2$，$\beta_1=\beta$）上类似地定义运算 $\beta_k$，其取值确定到模 $\Image\beta_{k-1}$：
\begin{equation}
\beta_k(z)=\frac{1}{p^k}\partial^*\tilde z\ (\mathrm{mod}\ p),\quad z\equiv\tilde z\ (\mathrm{mod}\ p),\ \ \tilde z\in C^q(X;\Z), \tag{6.12}
\end{equation}
\[
\beta_k:\Ker\beta_{k-1}\longrightarrow H^{q+1}(X;\Z/p)/\Image\beta_{k-1},
\]
\[
\Ker\beta_{k-1}\subset H^q(X;\Z/p).
\]
若对一切 $k\ge 1$ 都有 $\beta_k(z)=0$，则类 $z\in H^q(X;\Z/p)$ 是由一个整类约化得到的。类 $\beta_k(z)$ 总是来自 $p^k$ 阶的整类。
\end{example}

\begin{example}{2}
设 $k$ 是环或域，$z_j\in H^{n_j}(X;k)$（$j=1,2,3$）满足 $z_1z_2=0$，$z_2z_3=0$。对这种三元组，有一个称为 \emph{Massey 括号}的运算
\begin{equation}
\langle z_1,z_2,z_3\rangle=\theta(z_1,z_2,z_3)\subset H^{n_1+n_2+n_3-1}(X;k), \tag{6.13}
\end{equation}
它只定义到相差形如
\[
w_1z_3\pm z_1w_2,\qquad w_1\in H^{n_1+n_2-1}(X;k),\quad w_2\in H^{n_2+n_3-1}(X;k)
\]
的元素为止。Massey 括号的定义非常简单：在上链 $C^{n_i}(X;k)$ 上，它由
\[
\langle z_1,z_2,z_3\rangle=\partial^{-1}(z_1z_2)z_3\pm z_1\partial^{-1}(z_2z_3)
\]
给出。当 $k$ 是特征为 0 的域时，Massey 括号及其类似物，就复合而言，构成全部多值的、部分有定义的、自然的上同调运算的生成元。
\end{example}

现在回到关于 Steenrod 平方的讨论。由上面给出的分类（特别见 (6.11)）可以看出，Steenrod 代数 $\mathcal{A}_2$ 由形如 $Sq^{2^i}$（$i\ge 0$）的 Steenrod 平方乘性地生成。这一事实有下述很好的应用。

\begin{example*}
设 $K$ 是这样一个 CW 复形：
\[
H^j(K;\Z/2)\cong\Z/2\qquad\text{对 } j=0,n,2n,
\]
\[
H^j(K;\Z/2)=0\qquad\text{对 } j\ne 0,n,2n.
\]
记 $u_j$ 为 $H^j(K;\Z/2)$（$j=n,2n$）的生成元，并假定 $u_n^2=u_{2n}$。这样的 $K$ 的例子有实、复、四元数以及 Cayley 射影平面。
% [原文 p.95]
由此可知 $Sq^nu_n=u_{2n}$（由 Steenrod 平方的性质 (i)）。若 $n$ 不是 $2$ 的幂，这不可能，因为此时 $Sq^n$ 可分解为诸 $Sq^j$（$j<n$）的复合，而对这样的 $j$（$j=0$ 除外），由于假定 $H^{n+j}(K;\Z/2)=0$，我们有 $Sq^ju_n=0$。因此必须有 $n=2^s$。这一点 Hopf 在 20 世纪 40 年代不用上同调运算就已证明，Steenrod 在 40 年代末又用上同调运算证明了一次。

Adams 在 20 世纪 50 年代末进一步证明了：当 $n=2^i$（$i>3$）时，Steenrod 平方 $Sq^n$ 可非平凡地分解为部分运算的复合。特别地，由此可知 $Sq^nu_n=u_{2n}$ 仅当 $n=1,2,4,8$ 时成立。这一结果有一系列根本性的推论，特别是下述几条：
\begin{enumerate}[label=\arabic*.,leftmargin=2em]
\item 仅有的实可除代数（体）是 $\R^1,\R^2,\R^4,\R^8$。
\item 球面 $S^{n-1}$ 的稳定切丛平凡，仅当 $n=1,2,4,8$。
\item 群 $\pi_{2n-1}(S^n)$ 中存在 Hopf 不变量为 $1$ 的元素，仅当 $n=2,4,8$。\footnote{关于「Hopf 不变量」的定义，见第 4 章 §3。}
\item 存在 Hopf 型纤维化 $S^{2n-1}\to S^n$（即底空间与全空间都是球面的纤维化），仅当 $n=2,4,8$。\hfill$\square$
\end{enumerate}
\end{example*}

对 CW 复形 $K$ 构造 Steenrod 平方的基本想法如下。考虑对角
\[
\Delta(K)\subset K\times K,\qquad\Delta(K)=\{(x,x)\},
\]
并用 $\overline\Delta$ 表示 $K\times K$ 中包含 $\Delta$ 的那些胞腔构成的子复形。（注意 $\Delta$ 一般不是子复形。）设 $\Delta_1:K\to\overline\Delta\subset K\times K$ 是对角的一个胞腔逼近，即一个接近对角映射 $K\to\Delta$ 的胞腔映射，它把 $K$ 的每个胞腔 $\sigma^k$ 的像逼近地放在 $\overline\Delta(\sigma^k\times\sigma^k)$ 中。由此得到胞腔上链中的一个乘积运算
\[
C^*(K\times K)\cong C^*(K)\otimes C^*(K),
\]
它由
\[
z_iz_j=\Delta_1^*(z_i\otimes z_j)
\]
定义。这个运算把 $H^*(K)$ 变成环；它就是前面为单纯复形所定义的上同调乘积——见 §2。

考虑由
\[
\sigma(x,y)=(y,x),\qquad\sigma^2=1
\]
给出的对合 $\sigma:K\times K\to K\times K$。尽管 $\sigma\Delta=\Delta$，映射 $\sigma$ 一般并不（在平凡情形之外）保持 $\Delta_1$：$\sigma\Delta_1\ne\Delta_1$。这一不对称性正是 Steenrod 平方的来源。设 $w$ 是 $C^*(\Delta_1(K))$ 中的任一上闭链（为简明起见，以下只在 $\Z/2$ 上讨论），并用 $w$ 如下定义一列上闭链 $w_0,w_1,\dots$：
% [原文 p.96]
\[
w_0=w,\qquad \partial^*w_1=w_0+\sigma^*w_0,\qquad \partial^*(w_1+\sigma^*w_1)=0,
\]
\[
\partial^*w_2=w_1+\sigma^*w_1,\qquad \partial^*(w_2+\sigma^*w_2)=0,\ \dots,
\]
\[
\partial^*w_i=w_{i-1}+\sigma^*w_{i-1},\qquad \partial^*(w_i+\sigma^*w_i)=0,\ \dots,
\]
每个 $w_i$ 的支集都在 $\Delta_1(K)$ 中。对上闭链 $z\in C^n(K;\Z/2)$，令
\[
(z^2)_i=\Delta_1^*(w_i)\in C^{2n-i}(K;\Z/2),
\]
其中 $w_0$ 取为 $\Delta_1^*(z\otimes z)$；然后把 $z$ 的第 $(n-i)$ 个 Steenrod 平方定义为上闭链 $(z^2)_i\in H^{2n-i}(K;\Z/2)$ 的上同调类：
\[
Sq^{n-i}(z)=(z^2)_i.
\]
由这个定义立即可知：当 $i>0$ 时 $Sq^{n+i}(z)=0$，且 $Sq^n(z)=z^2$。其余性质的证明则较难；特别地，由上述构造并不能明显看出公理式地简单的性质 $Sq^0(z)=z$，即 $Sq^0=1$。另一方面，这个性质很容易由上面提到的 Serre 关于稳定 $\Z/2$-上同调运算的分类推出（据此，$j$ 次稳定运算构成的群同构于 $H^{j+n}(K(\Z/2,n);\Z/2)$）。因为 Eilenberg--MacLane 空间 $K(\Z/2,n)$ 的上同调在维数小于 $n$ 时平凡，而 $H^n(K(\Z/2,n);\Z/2)\cong\Z/2$，所以零次上同调运算只有一个，即乘以标量 $1\in\Z/2$。

% [Tutor Note: 递推式 $\partial^*w_i=w_{i-1}+\sigma^*w_{i-1}$ 的含义是：$w_{i-1}+\sigma^*w_{i-1}$ 恰是一个上边缘，故可解出 $w_i$；这一步之所以可行，全靠 $\sigma$ 是链映射且 $\sigma^2=1$。验证时可对 $i=1$ 逐个上链算一遍。]

上述 Steenrod 平方的构造在分次 \emph{Hopf 代数}的纯代数理论中有推论。设 $A$ 是 $\Z/p$ 上的分次代数，带单位元：
\[
A=\bigoplus_{n\ge0}A^n,\qquad A^0=\Z/p,
\]
乘法为
\[
\phi:A\otimes A\longrightarrow A,\qquad \phi(a\otimes b)=a\cdot b,\qquad A^n\cdot A^m\subset A^{n+m}.
\]
我们有增广同态 $\varepsilon:A\to\Z/p=A^0$，其中 $\Ker\varepsilon=\bigoplus_{n\ge1}A^n$。张量积（在 $\Z/p$ 上）
\[
A\otimes_{\Z/p}A,\qquad (A\otimes A)^j=\bigoplus_{q+m=j}A^q\otimes A^m
\]
是一个代数，其乘法由
\[
(a\otimes b)(a'\otimes b')=aa'\otimes bb'
\]
定义。

% [原文 p.97]
代数 $A\otimes A$ 有一个「扭」同态
\[
\sigma:A\otimes A\longrightarrow A\otimes A,\qquad \sigma(a\otimes b)=(-1)^{\deg a\deg b}b\otimes a.
\]
若存在一个\emph{对角同态}，即形如
\[
\Delta^*:A\longrightarrow A\otimes A,
\]
\[
\Delta^*(a)=a\otimes1+1\otimes a+\sum_j a'_j\otimes a''_j,\qquad \deg a'_j>0,\quad \deg a''_j>0
\]
的零次代数同态，则称代数 $A$ 为 \emph{Hopf 代数}。这里 $a\otimes b=(a\otimes1)\cdot(1\otimes b)$，而 $A\otimes1$ 与 $1\otimes A$ 在乘法下逐元素地交换到相差 $\pm1$。若对角同态 $\Delta^*$ 与由 $\sigma(a\otimes b)=(b\otimes a)$ 定义的对合算子 $\sigma:A\otimes A\to A\otimes A$ 交换（相差一个依赖于 $\deg a$ 的符号 $\pm$）：
\[
\Delta^*\sigma=\pm\sigma\Delta^*,
\]
则称 Hopf 代数 $A$ 是\emph{对称的}。

Milnor 在 20 世纪 50 年代末注意到，利用乘积上闭链上的作用公式（见 (6.10)），Steenrod 代数 $\mathcal{A}_p$ 可以自然地赋以 Hopf 代数结构。例如当 $p=2$ 时，令
\[
\Delta^*(Sq^i)=Sq^i\otimes1+1\otimes Sq^i+\sum_{j+k=i}Sq^j\otimes Sq^k.
\]
对任意 $\Z/p$ 上的对称 Hopf 代数 $A$，可以在相应的上同调代数 $\Ext_A^{*,*}(\Z/p,\Z/p)$（定义见下）上定义 Steenrod 运算的类似物：
\begin{equation}
St^i_p:\Ext_A^{m,n}(\Z/p,\Z/p)\longrightarrow\Ext_A^{m+i,pn}(\Z/p,\Z/p). \tag{6.14}
\end{equation}
例如，运算 $St^0_p$ 与由代数 $A$ 上的 Frobenius--Adams 算子 $\alpha:x\mapsto x^p$ 所决定的那个非平凡算子 $\alpha^*$ 一致。\footnote{当 $A$ 是 $\Z/p$ 上的 Hopf 代数时，它是有定义的。} 运算 $St^i_p$ 的性质与原来的 Steenrod 运算大致相同；而当 $p>2$ 时，$St^i_p\equiv0$ 仅对 $i\not\equiv0,1\pmod{p-1}$ 成立。因此，这里的非零运算 $St^i_p$ 比原来的 Steenrod 运算更多。

这些运算是在发展计算稳定同伦群所需的代数技巧的过程中发现的（这些计算用 Adams 谱序列进行：Novikov 在 20 世纪 50 年代末，Liulevicius 在 60 年代初）。\footnote{这些运算是 Novikov 为计算 Steenrod 代数 $\mathcal{A}_p$ 以及 $\mathcal{A}_p$ 上模的上同调而引入的。特别地，由此发现：对任意素数 $p$，球面的稳定同伦环中存在非平凡的挠元素 $x$ 使 $x^p\ne0$（1959）。Toda (1960) 用不同方法也证明了这一结果。Nishida 证明了如下定理：对球面同伦环 $\pi_*^s$ 的任一非平凡挠元素 $x$，存在 $n$ 使 $x^n=0$。因此环 $\pi_*^s$ 是局部幂零的，尽管不是强意义下的幂零。20 世纪 60 年代末，P. May 发表了一篇文章，用（略为更一般的范畴语言）建立了这些运算所需的全部基础结果。}\footnote{原书英译本译者注：对任意谱（Hopkins、Devinatz 与 Smith，20 世纪 80 年代中）已证明了一般的幂零定理。该定理最简单的形式是说：若「环谱」$X$（或具有乘法结构的「稳定空间」，例如上述稳定 Thom 空间）具有无挠上同调 $H^*(X;\Z)$，则 $\Tor\pi_*^s(X)$ 的每个元素都是幂零的。（细节及进一步的发展见 D. Ravenel, \emph{Nilpotence and Periodicity in Stable Homotopy Theory}, Annals of Mathematics Studies no.~128, Princeton, 1992.）}

% [原文 p.98]
在球面的 Adams 谱序列中，Steenrod 代数 $\mathcal{A}_p$ 的结构及其上同调代数
\begin{equation}
\Ext_{\mathcal{A}_p}^{*,*}(\Z/p,\Z/p) \tag{6.15}
\end{equation}
起着关键作用。按照 Adams 的方法，代数 (6.15) 的计算要用有限生成的对称 Hopf 代数对 Steenrod 代数 $\mathcal{A}_p$ 作「逼近」；此时上述 Steenrod 运算的类似物非常有用。当 $p>2$ 时，还必须同时考虑 Steenrod 代数 $\mathcal{A}_p$ 的双分次（Cartan 型）以及逼近它的有限 Hopf 代数。

把 Adams 谱序列用于计算任意 CW 复形 $K$ 的稳定同伦群时，人们把上同调代数 $H^*(K;\Z/p)=M$ 看作 Steenrod 代数上的分次左模。更一般地，考虑分次 Hopf 代数 $A$（在 $\Z/p$ 上）上的任意左模 $M$。模 $M$ 的上同调代数
\begin{equation}
\Ext_A^{*,*}(M,\Z/p) \tag{6.16}
\end{equation}
可以如下定义（这是 Cartan--Eilenberg 在 20 世纪 50 年代中叶的做法）。记 $A^+$ 为增广同态 $\varepsilon$ 的核 $\bigoplus_{n\ge1}A^n$，并考虑复形
\begin{equation}
M\leftarrow A^+\otimes M\leftarrow A^+\otimes A^+\otimes M\leftarrow\cdots\leftarrow A^+\otimes A^+\otimes\cdots\otimes A^+\otimes M\leftarrow\cdots, \tag{6.17}
\end{equation}
其中「$n$-链」具有形状
\[
\sigma^n=(a_1\otimes a_2\otimes\cdots\otimes a_n\otimes x)\in A^+\otimes A^+\otimes\cdots\otimes A^+\otimes M,
\]
\[
a_j\in A^{n_j},\quad n_j>0,\quad x\in M.
\]
$n$-链（或「单形」）$\sigma^n$ 的维数取为 $n$，并按它的次数分次，即令
\[
\deg\sigma^n=\sum_{j=1}^{n}n_j+n.
\]

% [原文 p.99]
边缘同态由
\[
\partial\sigma^n=\sum_{j=1}^{n-1}(a_1\otimes\cdots\otimes a_j\cdot a_{j+1}\otimes\cdots\otimes a_n\otimes x)+(-1)^n(a_1\otimes\cdots\otimes a_{n-1}\otimes a_n\cdot x)
\]
定义，并保持链的次数。复形 (6.17) 关于上述微分的同调就是所求的上同调代数 (6.16)。特别取 $M=\Z/p=H^*(S^0)$，便得到上面提到的上同调代数
\[
\Ext_{\mathcal{A}_p}^{*,*}(\Z/p,\Z/p).
\]
若模 $M$ 具有与 $\mathcal{A}_p$ 的对角相容的对角同态
\[
\Delta^*:M\longrightarrow M\otimes M,
\]
则上同调代数 $\Ext_{\mathcal{A}_p}^{*,*}(M,\Z/p)$ 获得分次 $\Z/p$-代数结构。

现在可以构造 \emph{Adams 谱序列}：它由群（实际上是 $\Z/p$-向量空间）$E_m=\bigoplus_{q,l}E_m^{q,l}$ 及微分（同态）
\[
d_m:E_m^{q,l}\longrightarrow E_m^{q+m,l+m-1},\qquad d_m^2=0
\]
构成，使得
\begin{enumerate}[label=(\roman*),leftmargin=2.2em]
\item $E_2^{q,l}\cong\Ext_{\mathcal{A}_p}^{q,l}(M,\Z/p)$，$E_{m+1}=H(E_m,d_m)$；
\item $\bigoplus\limits_{q-l=t}E_\infty^{q,l}\cong G\pi_t^s(K)$，其中 $G\pi$ 是「相伴群」。
\end{enumerate}
这里 $\pi_t^s(K)=\pi_{n+t}(\Sigma^nK)$，其中当 $j<n$ 时 $\pi_j(\Sigma^nK)=0$，$0\le t\le n-1$。相伴群 $G\pi$（关于滤过 $\pi=\pi^{(0)}\supset\pi^{(1)}\supset\cdots$）与环等的情形一样，是相继因子的直和：
\[
G\pi=\bigoplus_{j\ge0}\pi^{(j)}/\pi^{(j+1)}=\bigoplus_{j\ge0}G_j\pi.
\]
由 Adams 谱序列产生的滤过的所有成员之交（对给定的素数 $p$），恰好是 $K$ 的稳定同伦群中阶数与 $p$ 互素的有限阶元素之集。

当 $K=S^n$ 时（以及在另一些重要情形下），Adams 谱序列的各项是分次 $\Z/p$-代数 $E_m$，其微分 $d_m$ 满足 $d_m(uv)=d_m(u)v\pm ud_m(v)$（对分次元素 $u,v\in E_m$），
% [Tutor Note: Adams 谱序列的「次数」约定容易被 $q$、$l$、$t$ 三个指标绕晕：$q$ 是 Ext 的滤过次数（同调次数），$l$ 是内次数，收敛到 $\pi_t^s$ 时取 $t=q-l$；微分 $d_m$ 使 $q$ 增加 $m$、$l$ 增加 $m-1$，故 $q-l$ 恰好增加 $1$——这是理解全谱序列结构的钥匙。]
% [原文 p.100]
并且代数 $E_\infty$ 被附加到球面 $S^n$ 的稳定同伦群环上：
\[
\pi_*^s=\bigoplus_{k\ge0}\pi_k^s,\qquad \pi_k^s=\pi_{n+k}(S^n),\quad n>k+1,
\]
其乘法由复合 $S^{n+i+j}\xrightarrow{a}S^{n+i}\xrightarrow{b}S^n$ 诱导；通过把球面的稳定同伦群实现为「标架流形的配边环」（见第 4 章 §3），这个乘法对应于流形的直积所诱导的乘法。

对球面 $K=S^n$ 的情形，Adams 谱序列的第二项是代数
\[
E_2^{*,*}=\Ext_{\mathcal{A}_p}^{*,*}(\Z/p,\Z/p),
\]
其中 $\Z/p$ 的 $\mathcal{A}_p$-模结构由 $1\in\mathcal{A}_p$ 的恒等作用以及正次数元素对它的零作用决定。这里有一个元素
\[
h_0\in E_2^{1,1}=\Ext_{\mathcal{A}_p}^{1,1}(\Z/p,\Z/p),
\]
其性质是：乘以 $h_0$ 与乘以 $p$ 相伴。

\begin{example}{1}
由 $\mathcal{A}_p$-模 $M$ 的上同调的定义（见 (6.17)）可知
\begin{equation}
E_2^{0,k}=\Ext_{\mathcal{A}_p}^{0,k}(M,\Z/p)=\Hom_{\mathcal{A}_p}^k(M,\Z/p), \tag{6.18}
\end{equation}
其中只容许那些与 $\mathcal{A}_p$ 的作用交换的同态。$\mathcal{A}_p$ 在 $\Z/p$ 上的作用由
\[
a(u)=0\quad\text{若 }\deg a>0,
\]
\[
\lambda(u)=\lambda u\quad\text{若 }a=\lambda\in\Z/p
\]
给出，其中 $u\in\Z/p$。$\Hom_{\mathcal{A}_p}^k(M,\Z/p)$ 中的上标 $k$ 表示我们只容许在 $M^k$ 上为零的 $\Z/p$-线性形式（当然还要求它与 $\mathcal{A}_p$ 的作用交换）。当 $M=H^*(K;\Z/p)$ 时，零线 $E_2^{0,*}$ 有清楚的拓扑意义：任一映射 $S^k\xrightarrow{f}K$ 诱导 $\mathcal{A}_p$-模同态 $M=H^*(K;\Z/p)\xrightarrow{f^*}\Z/p$，从而决定一个元素 $f^*\in\Hom_{\mathcal{A}_p}^*(M,\Z/p)$；这个元素非零，保证映射 $f$ 是本质的（即不同伦于零）。反之，元素 $x\in\Hom_{\mathcal{A}_p}^*(M,\Z/p)$ 由某个映射 $S^n\to K$ 诱导，当且仅当 $x$ 在 Adams 谱序列的所有微分下都存活：
\[
d_m(x)=0,\qquad m\ge2.
\]

% [原文 p.101]
注意有许多本质映射 $S^k\xrightarrow{f}K$ 诱导平凡同态 $H^*(K;\Z/p)\xrightarrow{f^*}\Z/p$；例如球面稳定同伦群中正次数的非平凡元素，没有一个能由上同调环的相应同态「检测」出来。

\end{example}

\begin{example}{2}
如前所述，对球面 $S^n$ 有 $M=\Z/p$。此时 $\Ext_{\mathcal{A}_p}^{0,0}=\Z/p$，且当 $k\ne0$ 时 $\Ext_{\mathcal{A}_p}^{0,k}=0$。不难确定群 $\Ext_{\mathcal{A}_p}^{1,k}(\Z/p,\Z/p)$。为简明起见假定 $p=2$，简述其论证。由于代数 $\mathcal{A}_2$ 由元素 $Sq^{2^i}$ 生成，故存在非零元素
\[
h_i\in\Ext_{\mathcal{A}_2}^{1,2^i}(\Z/2,\Z/2),
\]
它与 $Sq^{2^i}$ 对偶。事实上这些元素穷尽了群 $\Ext_{\mathcal{A}_2}^{1,k}(\Z/2,\Z/2)$ 的全部非零元素。注意 $d_qh_i=0$（$q\ge2$）当且仅当球面的相应稳定同伦群中存在 Hopf 不变量为 $1$ 的元素（见第 4 章 §3）。Adams 在 20 世纪 50 年代末证明了
\[
d_2(h_i)=h_0h_{i-1}^2,
\]
其中当 $i\ge4$ 时 $h_0h_{i-1}^2\ne0$。对 $i=4$ 的情形不难验证 Adams 的公式：$d_2(h_4)=h_0h_3^2$。由于 $d_q(h_4)=0$（$q\ge3$：依「维数上的理由」，没有非平凡元素可作为靶），元素 $h_3$ 代表群 $\pi_7^s=\pi_{n+7}(S^n)$（$n\ge8$）中 Hopf 不变量为 $1$ 的非零元素 $h_3=\sigma$。由斜对称性我们有 $h_3\bar h_3=-\bar h_3h_3$，从而 $2\bar h_3^2=0$。由于乘以 $h_0$ 与乘以 $2$ 相伴，元素 $h_0h_3^2$ 在 Adams 谱序列的无穷项 $E_\infty$ 中必为平凡的，于是
\[
h_0h_3^2=dz\qquad\text{对某个 } z\in\Ext_{\mathcal{A}_2}^{1,16}(\Z/2,\Z/2),
\]
而由于 $h_4$ 是唯一的非零元素，必须有 $z=h_4$，这就建立了该公式。我们指出，2 线（2-line）$\Ext_{\mathcal{A}_2}^{2,*}(\Z/2,\Z/2)$ 由诸乘积
\begin{equation}
h_ih_j\in\Ext_{\mathcal{A}_2}^{1,2^i+2^j}(\Z/2,\Z/2), \tag{6.19}
\end{equation}
生成，其中 $h_ih_{i+1}=0$，而其它一切乘积都非平凡。乘积 $h_j^2$ 在流形论中特别重要，将于下一章讨论。\footnote{原书英译本译者注：元素 $h_j^2$（它们与「Arf 不变量问题」有关）是唯一一类人们尚不知道是否承载非平凡微分的双乘积 $h_ih_j$。已知 $h_j^2$ 对 $j=1,2,3,4,5$ 都是所有微分的闭链（M. Barratt、J. D. S. Jones 与 M. Mahowald，20 世纪 70 年代中）。元素 $h_ih_j$（$j\ge4$）非平凡，且当 $j\ge5$ 时其阶至少为 $4$（M. Mahowald, 1978）。关于球面 Adams 谱序列计算结果的一般综述见 D. Ravenel, \emph{Complex Cobordism and Stable Homotopy of Spheres}, Academic Press, 1986.}
\end{example}

% [Tutor Note: 原书把 (6.19) 印作 $\Ext^{1,2^i+2^j}$，但由 $h_i\in\Ext^{1,2^i}$ 相乘得 $h_ih_j\in\Ext^{2,2^i+2^j}$，次数必为 $2$；紧接的句子又称之为「2 线」，故上标 $1$ 系原书笔误，正确应为 $\Ext^{2,\,2^i+2^j}$。译文照原书印法保留。]

% [原文 p.102]
Adams 谱序列对光滑流形的配边环与协边环计算的应用，是 Milnor 与 Novikov 在 1960 年前后作出的。Adams 谱序列可以用超常（广义）上同调理论的术语自然地推广（Novikov，20 世纪 60 年代末），由此揭示出它与稳定同伦论的深刻联系；而计算稳定（及非稳定）同伦群的 Cartan--Serre 方法则不容许这样的推广。

在有理数域上，Adams 谱序列收敛于稳定有理同伦群 $\pi_t^s(K)\otimes\Q=\pi_{n+t}(\Sigma^nK)$（$n$ 充分大），并归结为一项平凡的例行步骤，给出 Serre 早先（20 世纪 50 年代初）的结果：Hurewicz 同态事实上是同构：
\begin{equation}
\pi_t^s(K)\otimes\Q\xrightarrow{\cong}H_t(K;\Q). \tag{6.20}
\end{equation}
如前所述，有理上同调理论的运算代数 $\mathcal{A}_\Q$ 只由数乘组成。

也如前所述，对素数 $p$ 把 Adams 谱序列用于同伦论的具体问题时，需要研究 $\Z/p$ 上的各种 Hopf 代数：
\[
B=\bigoplus_{n\ge0}B^n,\qquad B^0=\Z/p,
\]
它们带有对称对角 $\Delta:B\to B\otimes B$（见上）。这种代数（态射保持分次与 Hopf 结构）的范畴以及 Hopf 代数上模的范畴，构成一个有趣的代数模型\footnote{特别地，在这些范畴中，对 Hopf 代数 $A$ 与 $A$-模 $R$ 定义了上同调代数 $H^{*,*}(A;R)$。}，它模仿 CW 复形的同伦范畴的某些性质。例如，底空间单连通的纤维丛在代数上由 Hopf 代数的一个适当满射 $f:A\to B$ 来模拟（其中 $A$ 是「全空间」，$B$ 是「底空间」），并具有下述特殊性质：应当存在 Hopf 子代数 $C\subset A$（「纤维」）：
\[
C=\bigoplus_{n\ge0}C^n\subset A=\bigoplus_{n\ge0}A^n,\qquad C^0=A^0=\Z/p,
\]
其增广理想
\[
C^+=\bigoplus_{n\ge1}C^n\subset A^+=\bigoplus_{n\ge1}A^n,
\]
使得
\begin{enumerate}[label=(\roman*),leftmargin=2.2em]
\item $C$ 是 $A$ 的中心子代数（对一切 $c\in C$、$a\in A$ 有 $ac=ca$），且 $\Ker f$ 具有形状 $\Ker f=AC^+=C^+A$；
\end{enumerate}
% CHUNK ch03g DONE
